%========================================================
% ABSTRACT IN ENGLISH
%========================================================

As quantum technologies progress toward distributed architectures, optical quantum links are becoming 
increasingly important for both interconnecting remote quantum processing nodes and enabling secure quantum communication. 
Polarization-encoded photonic qubits are particularly suited to fiber-based transmission, since polarization can be prepared, 
manipulated, and measured using well-established optical components.

Among the applications considered in this thesis, particular attention is devoted to \\entanglement-based quantum key distribution, 
with the BBM92 protocol providing the main operational reference. 
In this context, the preservation of high-quality entanglement is directly related to the suitability of the distributed state for secure quantum communication. 
The \textit{Quantum Bit Error Rate} (QBER) provides a direct protocol-level criterion for assessing whether the measured correlations 
are compatible with secure key extraction, while the \textit{Secure Key Rate} (SKR) represents the corresponding operational throughput once experimental 
parameters and security quality are taken into account.

This thesis develops a progressive and unified framework for studying the propagation of polarization-entangled photon pairs through optical fibers, 
connecting physical descriptions commonly treated separately at different levels of approximation. The main objective is to relate fiber and source parameters 
to the quality of the distributed entanglement and to the operational requirements of BBM92-based communication. 
To this purpose, the framework connects the physical parameters of the optical link with state-quality metrics and with protocol-relevant quantities derived from the polarization correlations, particularly the QBER.
This provides a practical procedure for identifying operating regions of the physical parameter space in which the predicted polarization correlations satisfy an adopted BBM92 QBER criterion, 
without assuming a universal concurrence threshold.

The modeling hierarchy starts from the residual polarization evolution after partial compensation. A fixed relative phase first distinguishes 
coherent changes in measured correlations from genuine entanglement degradation. The model is then extended to unresolved phase fluctuations through the coherence 
parameter \(\mu\), and to independent depolarizing effects in the two fiber arms through depolarizing parameters \(p_A\) and \(p_B\). 
Finally, these effective descriptions are connected to polarization-mode dispersion (PMD), where finite-bandwidth photons undergo frequency-dependent birefringent 
transformations and polarization mixedness emerges after tracing over unresolved spectral degrees of freedom.

The framework is implemented in a modular MATLAB simulation environment that integrates quantum-state preparation, 
fiber-channel modeling, statistical averaging, spectral tracing, and polarization-correlation analysis. 
Analytical and numerical results are used to study how source properties, fiber parameters, PMD, and 
compensation conditions affect the surviving entanglement.

The theoretical and numerical analysis is complemented by a preliminary laboratory study of a simplified two-arm, BBM92-oriented polarization platform. 
The experimental activity is designed to characterize, under controlled conditions, the polarization transformations 
and degradation mechanisms introduced by the optical link and by selected fiber components.
Although the present setup does not directly implement a complete entanglement-based QKD protocol, 
it provides a first experimental step toward assessing the practical requirements of fiber-based entanglement distribution and toward establishing the connection between 
fiber-induced polarization degradation and, in a future complete BBM92 implementation, experimentally measurable quantities such as the QBER and the Secure Key Rate.

\medskip
\noindent
\textbf{Keywords:}
Polarization Entanglement,
Optical Fibers,
BBM92,
Quantum Key Distribution,
Secure Key Rate,
Quantum Channels,
Decoherence,
Depolarization.
